% Repeat the complete overview with the QPIP_1 enclosure already visible.
\begin{frame}{Protocol Overview}
  \centering
  \begin{tikzpicture}[
      protocolbox/.style={draw=none,fill=CVQCPale,rounded corners=4mm},
      modellabel/.style={inner sep=0pt,
        font=\small\bfseries,text=CVQCInk}
    ]
    \path[use as bounding box] (-6.6,-3.20) rectangle (6.6,3.05);
    \path[protocolbox] (-4.75,-3.15) rectangle (4.75,2.15);
    \node[modellabel,anchor=north west] at (-4.45,1.95)
      {$\mathsf{QPIP}_1$};
    \protocolflowchart
  \end{tikzpicture}

  \note{
    \footnotesize
    Return to the complete overview, with the QPIP1 part already enclosed.
    We have explained the central mechanism: authenticated distillation
    removes the plaintext-level attacks, leaving ciphertext-level attacks
    on the surviving encodings. Those are the form that authenticated
    evaluation is designed to handle. The full paper combines this with
    the preparation stage's joint bad-support bound, recursive counting,
    the other magic-state routines, and verified evaluation to obtain
    the negligible-error statistical QPIP1 guarantee. The local
    one-bad-input lemma alone is not the entire QPIP1 proof.
    For the rest of the talk, focus on the remote-preparation box above
    the enclosure. It replaces the weak verifier's single-qubit messages
    with classical interaction. We still need an input-quality and
    factorizability guarantee strong enough to feed the same construction,
    including the handling of complex conjugation. Leave the underlying
    classical-client QFHE implementation in the paper; introduce the
    HN message interface next.
  }
\end{frame}
